% ===========================================================================
%  Annexe — Glossaire des notions
% ===========================================================================
\chapter{Glossaire des notions}
\markboth{Glossaire}{Glossaire}

Un mot mal compris coûte plus cher qu'une formule oubliée. La formule, on la
retrouve dans le formulaire ; le mot, lui, commande toute la lecture de
l'énoncé. « Escompte », « barycentre », « équiprobable », « amortissement » :
ces termes ont, en mathématiques comme en économie, un sens étroit, et c'est
ce sens-là que le correcteur attend.

Les pages qui suivent rassemblent les termes de l'année dans l'ordre
alphabétique. Chaque définition est courte à dessein : elle rappelle ce que
le mot désigne, non comment on s'en sert. Le numéro du chapitre indiqué à
droite renvoie à l'endroit où la notion est construite, démontrée et mise au
travail ; quelques entrées renvoient à l'annexe de mathématiques financières
plutôt qu'à un chapitre numéroté.

% ---------------------------------------------------------------------------
% Mise en forme locale
\newcommand{\glolettre}[1]{%
  \par\vspace{1.15em}\noindent
  \begingroup\sffamily\bfseries\large\color{bmtAccent}#1\endgroup
  \hspace{0.5em}%
  {\color{bmtGrisClair}\leaders\hrule height 1pt\hfill\kern0pt}%
  \par\vspace{0.35em}}

\newcommand{\glo}[2]{%
  \par\vspace{0.5em}\noindent
  {\sffamily\bfseries\small\color{bmtAccent}#1}%
  \hspace{0.45em}{\sffamily\footnotesize\color{bmtGris}ch.~#2}%
  \par\vspace{0.05em}\noindent\ignorespaces}

% ---------------------------------------------------------------------------
\glolettre{A}

\glo{Affixe}{9}
Nombre complexe associé à un point ou à un vecteur du plan : le point
$M(a\,;b)$ a pour affixe $z = a+ib$. C'est le pont qui fait passer d'une
figure à un calcul.

\glo{Amortissement}{annexe}
Part d'une annuité qui rembourse le capital emprunté, par opposition à
l'intérêt, qui rémunère le prêteur. Dans un tableau d'amortissement, la
somme des amortissements successifs reconstitue le capital initial.

\glo{Annuité}{annexe}
Versement périodique — le plus souvent annuel — destiné à rembourser un
emprunt ou à constituer un capital. Une annuité constante se décompose
chaque période différemment entre intérêt et amortissement.

\glo{Antidéplacement}{11}
Isométrie du plan qui renverse l'orientation : les symétries axiales et les
symétries glissées, et elles seules.

\glo{Argument}{9}
Mesure de l'angle que fait avec l'axe des abscisses le vecteur d'affixe $z$
non nulle. Il est défini à $2k\pi$ près, ce qui oblige à préciser chaque fois
l'intervalle choisi.

\glo{Asymptote}{3}
Droite dont la courbe se rapproche indéfiniment. Verticale quand la fonction
explose en un point, horizontale ou oblique quand la différence entre la
courbe et la droite tend vers zéro à l'infini.

% ---------------------------------------------------------------------------
\glolettre{B}

\glo{Barycentre}{10}
Point d'équilibre d'un système de points affectés de coefficients. Il n'existe
que si la somme des coefficients n'est pas nulle, et il ne change pas quand on
multiplie tous les coefficients par un même réel non nul.

\glo{Bijection}{1, 5}
Fonction qui, à chaque valeur de l'ensemble d'arrivée, fait correspondre un
antécédent et un seul. Une fonction continue et strictement monotone sur un
intervalle en réalise toujours une.

\glo{Binomiale (loi)}{14}
Loi de probabilité du nombre de succès obtenus en répétant $n$ fois, de façon
indépendante, une même épreuve à deux issues.

% ---------------------------------------------------------------------------
\glolettre{C}

\glo{Combinaison}{13}
Sous-ensemble de $p$ éléments pris parmi $n$, sans considération d'ordre. Le
nombre de ces sous-ensembles est noté $\combi{p}{n}$.

\glo{Conjugué}{9}
Nombre complexe $\overline z = a-ib$ associé à $z = a+ib$. Géométriquement,
c'est le symétrique de $z$ par rapport à l'axe des réels.

\glo{Continuité}{1}
Propriété d'une fonction dont la limite en un point est égale à sa valeur en
ce point. Sur un intervalle, elle traduit une courbe que l'on trace sans lever
le crayon.

\glo{Corrélation (coefficient de)}{15}
Nombre compris entre $-1$ et $1$ qui mesure la force d'une liaison linéaire
entre deux séries statistiques. Plus il est proche de $1$ en valeur absolue,
plus les points du nuage s'alignent sur une droite.

\glo{Convexité}{2}
Une fonction est convexe sur un intervalle si sa courbe y reste au-dessus de
chacune de ses tangentes ; cela équivaut à une dérivée seconde positive.

\glo{Covariance}{15}
Moyenne des produits des écarts à la moyenne de deux séries statistiques
appariées. Son signe indique le sens de la liaison ; sa valeur entre dans le
calcul du coefficient de corrélation et de la pente de la droite de
régression.

% ---------------------------------------------------------------------------
\glolettre{D}

\glo{Déplacement}{11}
Isométrie du plan qui conserve l'orientation : translations et rotations, et
elles seules.

\glo{Dérivabilité}{2}
Existence d'une limite finie du taux d'accroissement en un point. Cette limite
est le nombre dérivé, et c'est le coefficient directeur de la tangente.

% ---------------------------------------------------------------------------
\glolettre{E}

\glo{Écart type}{14}
Racine carrée de la variance. Il mesure la dispersion des valeurs autour de
l'espérance, dans la même unité que la variable.

\glo{Équiprobabilité}{13}
Situation où toutes les issues d'une expérience ont la même probabilité. La
probabilité d'un événement se réduit alors au quotient du nombre de cas
favorables par le nombre de cas possibles.

\glo{Escompte}{annexe}
Réduction accordée lorsqu'un effet de commerce est négocié avant son
échéance : le porteur cède l'effet contre son montant diminué de l'intérêt
correspondant au temps restant à courir.

\glo{Espérance}{14}
Moyenne des valeurs d'une variable aléatoire, chacune pesée par sa
probabilité. C'est le gain moyen que l'on obtiendrait en répétant très
longtemps l'expérience.

\glo{Événement}{13}
Partie de l'univers, c'est-à-dire ensemble d'issues. Il est réalisé dès que
l'issue observée lui appartient.

\glo{Exponentielle}{6}
Fonction réciproque du logarithme népérien, notée $\exp$ ou $x \mapsto e^{x}$.
Elle est sa propre dérivée, ce qui explique sa présence dans les calculs
d'intérêt composé.

% ---------------------------------------------------------------------------
\glolettre{F}

\glo{Fonction réciproque}{5, 6}
Fonction qui défait ce que fait une bijection : si $f(a) = b$, alors
$f^{-1}(b) = a$. Les deux courbes sont symétriques par rapport à la droite
d'équation $y = x$.

% ---------------------------------------------------------------------------
\glolettre{H}

\glo{Homothétie}{11, 12}
Transformation qui, à partir d'un centre $\Omega$ et d'un rapport $k$ non nul,
envoie $M$ sur $M'$ tel que $\vect{\Omega M'} = k\,\vect{\Omega M}$. Elle
multiplie les longueurs par $\abs k$ et conserve les angles.

% ---------------------------------------------------------------------------
\glolettre{I}

\glo{Indépendance}{13, 14}
Deux événements sont indépendants lorsque la réalisation de l'un ne change pas
la probabilité de l'autre, ce qui se traduit par
$\proba{A \cap B} = \proba A \times \proba B$.

\glo{Intégrale}{4}
Pour une fonction continue et positive, aire du domaine compris entre la
courbe, l'axe des abscisses et deux droites verticales. Dans le cas général,
c'est la variation d'une primitive entre les deux bornes.

\glo{Intérêt composé}{annexe}
Mode de calcul dans lequel l'intérêt de chaque période s'ajoute au capital et
produit lui-même un intérêt les périodes suivantes : le capital croît alors
de façon exponentielle.

\glo{Intérêt simple}{annexe}
Mode de calcul dans lequel l'intérêt, proportionnel au capital, au taux et à
la durée, ne s'ajoute pas au capital en cours de placement : il croît de
façon linéaire.

\glo{Isométrie}{11}
Transformation du plan qui conserve les distances. Elle conserve donc aussi
les angles géométriques, les milieux et les aires.

% ---------------------------------------------------------------------------
\glolettre{L}

\glo{Limite}{1}
Valeur dont $f(x)$ se rapproche autant qu'on veut lorsque $x$ s'approche d'un
point ou tend vers l'infini. Elle peut être finie, infinie, ou ne pas exister.

\glo{Logarithme népérien}{5}
Unique primitive de $x \mapsto \frac1x$ sur $\Rps$ qui s'annule en $1$. Il
transforme les produits en sommes, ce qui est la raison de son invention.

% ---------------------------------------------------------------------------
\glolettre{M}

\glo{Moindres carrés (méthode des)}{15}
Procédé qui détermine la droite de régression en rendant minimale la somme
des carrés des écarts verticaux entre les points du nuage et la droite.

\glo{Module}{9}
Distance de l'origine au point d'affixe $z$ : $\abs z = \sqrt{a^{2}+b^{2}}$
pour $z = a+ib$. Le module d'un produit est le produit des modules.

\glo{Monotone}{2, 3}
Se dit d'une fonction croissante sur tout un intervalle, ou décroissante sur
tout un intervalle. La stricte monotonie interdit les paliers.

% ---------------------------------------------------------------------------
\glolettre{N}

\glo{Nuage de points}{15}
Représentation graphique d'une série statistique à deux variables : chaque
individu y est un point de coordonnées $(x_{i}\,;y_{i})$.

% ---------------------------------------------------------------------------
\glolettre{P}

\glo{Primitive}{4}
Fonction dont la dérivée est la fonction de départ. Deux primitives d'une même
fonction sur un intervalle ne diffèrent que d'une constante.

\glo{Probabilité conditionnelle}{13}
Probabilité d'un événement $B$ sachant qu'un événement $A$ de probabilité non
nulle est réalisé, notée $\probacond BA$. Elle se lit directement sur une
branche d'arbre pondéré.

\glo{Point moyen}{15}
Point $G\left(\bar x\,;\bar y\right)$ dont les coordonnées sont les moyennes
des deux séries. La droite de régression passe toujours par ce point.

% ---------------------------------------------------------------------------
\glolettre{R}

\glo{Régression (droite de)}{15}
Droite d'équation $y = ax+b$ qui résume au mieux, au sens des moindres
carrés, la liaison entre deux séries statistiques ; elle sert à estimer $y$
pour une valeur de $x$ non observée.

\glo{Rotation}{11}
Isométrie déterminée par un centre et un angle. Elle conserve les distances et
l'orientation, et ne laisse fixe que son centre — sauf si l'angle est nul.

% ---------------------------------------------------------------------------
\glolettre{S}

\glo{Similitude}{12}
Transformation qui multiplie toutes les distances par un même rapport $k>0$.
Directe si elle conserve l'orientation, indirecte sinon ; les isométries sont
les similitudes de rapport $1$.

\glo{Suite arithmétique}{7}
Suite dont on passe d'un terme au suivant en ajoutant toujours le même nombre,
appelé raison.

\glo{Suite géométrique}{7}
Suite dont on passe d'un terme au suivant en multipliant toujours par le même
nombre non nul, appelé raison.

\glo{Symétrie glissée}{11}
Composée d'une symétrie axiale et d'une translation de vecteur dirigeant
l'axe. C'est l'antidéplacement sans point fixe.

\glo{Système linéaire}{8}
Ensemble d'équations du premier degré à résoudre simultanément. La méthode
du pivot de Gauss l'échelonne, équation après équation, jusqu'à ce qu'une
résolution en cascade devienne possible.

% ---------------------------------------------------------------------------
\glolettre{T}

\glo{Tangente}{2}
Droite qui approche la courbe au plus près au voisinage d'un point. Son
coefficient directeur est le nombre dérivé en ce point.

\glo{Tableau d'amortissement}{annexe}
Tableau qui détaille, période par période, l'intérêt, l'amortissement et le
capital restant dû d'un emprunt, jusqu'à extinction complète de la dette.

\glo{Théorème des valeurs intermédiaires}{1}
Une fonction continue sur un intervalle prend toute valeur comprise entre deux
de ses valeurs. Avec la stricte monotonie, la solution obtenue est unique.

\glo{Translation}{11}
Isométrie qui déplace tous les points d'un même vecteur. Elle n'a aucun point
fixe, sauf si le vecteur est nul.

% ---------------------------------------------------------------------------
\glolettre{U}

\glo{Univers}{13}
Ensemble de toutes les issues possibles d'une expérience aléatoire. Le choisir
correctement est la première décision d'un exercice de probabilités, et
souvent la seule difficile.

% ---------------------------------------------------------------------------
\glolettre{V}

\glo{Valeur acquise}{annexe}
Montant obtenu à une date future par un capital placé, intérêts compris —
qu'il s'agisse d'un placement unique ou d'une suite d'annuités.

\glo{Valeur actuelle}{annexe}
Montant, à une date antérieure, équivalent à un capital ou à une suite de
versements futurs, compte tenu du taux d'actualisation. C'est l'opération
inverse de la capitalisation.

\glo{Valeur nominale}{annexe}
Montant inscrit sur un effet de commerce, dû à son échéance, avant toute
déduction d'escompte.

\glo{Variable aléatoire}{14}
Fonction qui associe un nombre à chaque issue d'une expérience aléatoire. Sa
loi donne, pour chaque valeur prise, la probabilité correspondante.

\glo{Variance}{14}
Moyenne des carrés des écarts à l'espérance. Elle mesure la dispersion, mais
dans une unité au carré, d'où le passage à l'écart type.
